\section{Experimental Design}
\label{sec:design}

This section describes the comparisons reported in the main results. Evaluation
metrics, threshold selection, and batch selection are defined in
Section~\ref{sec:problem}.

\subsection{Main Benchmark}
\label{sec:design:main}

The primary comparison evaluates Kalman-BCE and Kalman-LSTM-Spec across all
three systems at five patient counts (100, 200, 300, 400, 500), each with ten
seeds. This produces 60~runs per patient count:
\(3\;\text{systems} \times 2\;\text{methods} \times 10\;\text{seeds}\).

Table~\ref{tab:benchmark_500} reports the 500-patient results: mean and
standard deviation of detection time, early-warning AUC, and false-positive rate
across seeds, listed per system and method. Table~\ref{tab:benchmark_sweep}
extends this to all five patient counts.

Figure~\ref{fig:patient_sweep} plots detection time and early-warning AUC as a
function of patient count, with separate panels per system. This tests whether
performance changes with training set size and whether the two variants respond
differently to that change.

\subsection{Trajectory-Level Behavior}
\label{sec:design:trajectory}

Figure~\ref{fig:trajectory_panels} shows representative test trajectories from
each system alongside the corresponding observer score sequences. The panels
display how scores rise in the pre-bifurcation region for signal trajectories
and stay flat for null trajectories. A method with high detection time should
show early score rises in these panels; a method with low false-positive rate
should produce near-zero scores on null sequences.

\subsection{Training Behavior}
\label{sec:design:training}

Figure~\ref{fig:training_curves} plots training and validation loss curves for
both variants at the 500-patient setting. These curves confirm that training
converges within the allotted 50~epochs and show whether early stopping
activates at a consistent epoch or varies across seeds.

\subsection{Supplementary and Exploratory Comparisons}
\label{sec:design:supplementary}

Several comparisons are excluded from the main tables. Classical early-warning
indicators (lag-1 autocorrelation, rolling variance) were evaluated during
development but use a different scoring pipeline and are not included in the
finalized benchmark. Ablation experiments---EWS-augmented input features,
auxiliary spectral heads, null-training configurations---were run with fewer
seeds or on a subset of systems. These results are retained in the experiment
directory for reference but are not mixed into the main claims, which report
only the two observer variants across the complete seed and system grid.
